\sheet{5}{Nyquist}{Chapter~5 (Sampling Theorem, Aliasing, Reconstruction)}

\begin{goals}
\begin{itemize}
  \item Predict alias frequencies with the folding rule and the folding
    diagram; determine minimum sampling rates, including band-pass
    (under-)sampling.
  \item Observe aliasing of the knock resonances on real (simulated)
    microcontrollers and verify the predictions quantitatively.
  \item Compare zero-order hold, linear and sinc reconstruction.
  \item Design an anti-aliasing filter: trade-off between filter order,
    sampling rate and passband loss.
\end{itemize}
\end{goals}

\section*{Background}

Sampling with $f_s=1/T_s$ makes all frequencies $f+kf_s$, $k\in\mathbb{Z}$,
indistinguishable. A real tone of frequency $f$ therefore appears at the
\emph{alias frequency}
\[
f_a=\bigl|\,f-k f_s\bigr|,\qquad k=\operatorname{round}(f/f_s),
\qquad 0\le f_a\le f_s/2 .
\]
A band-limited signal ($X(f)=0$ for $|f|\ge B$) is uniquely determined by its
samples if $f_s>2B$ and can be reconstructed by sinc interpolation
$x(t)=\sum_n x[n]\,\operatorname{sinc}\bigl((t-nT_s)/T_s\bigr)$,
$\operatorname{sinc}(u)=\sin(\pi u)/(\pi u)$. Everything above $f_s/2$ has to be
removed \emph{before} sampling by an analog anti-aliasing (AA) filter.

In our virtual engine the knock sensor has resonances at \SI{6.5}{kHz} and
\SI{10.5}{kHz} ($\tau=\SI{0.8}{ms}$) and valve impacts at \SI{8}{kHz}
($\tau=\SI{0.4}{ms}$). The \texttt{engine-sim} chip updates its analog outputs
every \SI{20}{\micro s} (attribute \texttt{fs\_khz = 50}) and holds the value in
between: the chip itself is a sampled system with zero-order hold.

%-------------------------------------------------------------------------
\begin{exercise}{Aliasing by hand}{\Paper}
\begin{tasks}
  \item Compute the alias frequencies of the components \SI{6.5}{kHz},
    \SI{8}{kHz}, \SI{10.5}{kHz} and of a (hypothetical) third knock mode at
    \SI{15}{kHz} for $f_s=5$, $8.33$, $16.67$, $25$ and \SI{50}{kHz}.
    (\SI{8.33}{kHz} is close to the maximum rate of the Uno ADC.)
  \item Draw the folding diagram ($f_a$ over $f$, \SIrange{0}{16}{kHz}) for
    $f_s=\SI{8.33}{kHz}$ and mark the four components. Which two components
    can no longer be separated after sampling?
  \item What is the minimum sampling rate for a knock signal with components
    up to \SI{15}{kHz}? A 4th-order Butterworth AA filter with
    $f_c=\SI{15}{kHz}$ ($|H|=1/\sqrt{1+(f/f_c)^8}$) is used. Which $f_s$ is
    needed so that every frequency that folds into \SIrange{0}{15}{kHz} is
    attenuated by at least \SI{40}{dB}?
  \item Band-pass sampling: an analog band-pass passes only the first knock
    mode, \SIrange{6}{7}{kHz}. Which sampling rates are allowed
    ($2f_H/k\le f_s\le 2f_L/(k-1)$, $k=1,2,\dots$)? Show that the Uno at
    \SI{5}{kHz} could legally capture this band. Where does the band appear,
    and is it spectrally inverted?
\end{tasks}
\end{exercise}

\begin{solution}
a) Alias frequencies in kHz ($k$ in brackets):
\begin{center}\small
\begin{tabular}{@{}lccccc@{}}\toprule
$f$ & $f_s=5$ & $8.33$ & $16.67$ & $25$ & $50$\\\midrule
6.5 (knock 1)  & 1.5 (1) & 1.833 (1) & 6.5 (0) & 6.5 & 6.5\\
8.0 (valve)    & 2.0 (2) & 0.333 (1) & 8.0 (0) & 8.0 & 8.0\\
10.5 (knock 2) & 0.5 (2) & 2.167 (1) & 6.167 (1) & 10.5 & 10.5\\
15.0 (knock 3) & 0 (3)   & 1.667 (2) & 1.667 (1) & 10.0 (1) & 15.0\\\bottomrule
\end{tabular}
\end{center}
Notable cases: at \SI{5}{kHz} the \SI{15}{kHz} mode becomes DC; at
\SI{16.67}{kHz} the second knock mode lands at \SI{6.17}{kHz}, right next to
the first mode -- a \SI{6.5}{kHz} band-pass would see it; at \SI{25}{kHz} a
\SI{15}{kHz} mode would mimic \SI{10}{kHz}.

b) Folding diagram for $f_s=\SI{8.33}{kHz}$ (zig-zag with corners at
multiples of $f_s/2$):
\begin{center}
\begin{tikzpicture}
\begin{axis}[width=10cm,height=4.6cm,xmin=0,xmax=16.5,ymin=0,ymax=4.6,
  xlabel={$f$ [kHz]},ylabel={$f_a$ [kHz]},xtick={0,4.17,8.33,12.5,16.67},
  xticklabels={0,$\frac{f_s}{2}$,$f_s$,$\frac{3f_s}{2}$,$2f_s$},
  ytick={0,1,2,3,4.17},yticklabels={0,1,2,3,$\frac{f_s}{2}$},
  grid=major,grid style={gray!25},font=\small]
\addplot[thick,dspblue] coordinates {(0,0) (4.1667,4.1667) (8.3333,0) (12.5,4.1667) (16.6667,0)};
\addplot[only marks,mark=*,red!70!black] coordinates {(6.5,1.8333) (8,0.3333) (10.5,2.1667) (15,1.6667)};
\node[font=\scriptsize,anchor=south east] at (axis cs:6.5,1.9) {6.5};
\node[font=\scriptsize,anchor=west] at (axis cs:8.1,0.35) {8.0};
\node[font=\scriptsize,anchor=south east] at (axis cs:10.5,2.2) {10.5};
\node[font=\scriptsize,anchor=south west] at (axis cs:15,1.7) {15};
\end{axis}
\end{tikzpicture}
\end{center}
The two knock modes end up at \SI{1.83}{kHz} and \SI{2.17}{kHz}: only
\SI{333}{Hz} apart, which is less than the frequency resolution of a knock
burst of a few ms (Sheet~8). Mode 3 (\SI{1.67}{kHz}) falls into the same
region; the valve impact becomes an almost-DC transient.

c) Nyquist: $f_s>\SI{30}{kHz}$. With an AA filter, a frequency $f$ folds into
$[0,15]$\,kHz if $f\ge f_s-\SI{15}{kHz}$; the weakest attenuation occurs at
$f=f_s-15$. $40$\,dB requires $(f/f_c)^8\ge10^4$, i.e.\
$f\ge f_c\cdot10^{0.5}=\SI{47.4}{kHz}$, so $f_s\ge\SI{62.4}{kHz}$ (in
practice \SI{64}{kHz}). The theoretical minimum of \SI{30}{kHz} would need an
ideal brick-wall filter.

d) $f_L=6$, $f_H=7$, $B=1$\,kHz, $k\le\lfloor f_H/B\rfloor=7$:
\begin{center}\small
\begin{tabular}{@{}cccccccc@{}}\toprule
$k$ & 1 & 2 & 3 & 4 & 5 & 6 & 7\\\midrule
$f_s$ [kHz] & $\ge14$ & 7--12 & 4.67--6 & 3.5--4 & 2.8--3 & 2.33--2.4 & 2\\\bottomrule
\end{tabular}
\end{center}
$f_s=\SI{5}{kHz}$ lies in the $k=3$ zone. The band lies in the third Nyquist
zone $[5,7.5]$\,kHz (odd zone, not inverted) and appears at $f-f_s$, i.e.\ at
\SIrange{1}{2}{kHz}. Condition: the analog band-pass must suppress the valve
(\SI{8}{kHz}$\to$\SI{2}{kHz}) and mode 2 (\SI{10.5}{kHz}$\to$\SI{0.5}{kHz}),
which fold into the same \SIrange{0}{2.5}{kHz} range. Band-pass sampling
replaces a high $f_s$ by a demanding analog filter.
\end{solution}

%-------------------------------------------------------------------------
\begin{exercise}{Aliasing made visible}{\Wokwi}
Two projects record the \texttt{KNOCK} output of the engine-sim chip (knock
probability \SI{100}{\%}, \SI{3000}{rpm}) together with the ground-truth
outputs \texttt{KNK} (HIGH during a knock burst) and \texttt{TDC}, and print
blocks of CSV lines \texttt{n,adc,knk,tdc}:
\codefile{wokwi/sheet05\_alias\_uno/} (Uno, Timer1 ISR, 512 samples per
block) and \codefile{wokwi/sheet05\_alias\_esp32/} (ESP32, deadline loop,
2500 samples per block). The JSLinux program
\codefile{jslinux/sheet05/alias\_scan.c} cuts out every knock burst (from the
rising KNK edge, $4\tau$) and every valve burst (TDC $+90^\circ$, $5\tau$),
averages the spectra $|X(f)|^2$ of the segments and lists the peaks.
\begin{tasks}
  \item Uno: the sampling period is \texttt{TS\_US = DIV}$\times$\SI{20}{\micro
    s} (default \texttt{DIV = 10}, \SI{5}{kHz}). Compute \texttt{OCR1A}
    (prescaler 8) and complete the ISR (store KNK in bit 14, TDC in bit 15).
    Why does the sketch disable the Timer0 interrupt during acquisition?
  \item Copy at least three blocks from the Serial Monitor into
    \texttt{uno.csv}, implement \texttt{alias\_freq()} and the start index
    of the valve burst in \texttt{alias\_scan.c} and run
    \texttt{./alias\_scan uno.csv 5000 10}. Compare the spectral peaks with
    Exercise~5.1. Look at one knock burst in the CSV: how many oscillation
    periods do you see, and of which frequency? Repeat with \texttt{DIV = 6}
    (\SI{8.33}{kHz}). Check with the logic analyzer that the sample marker
    (D0) is still regular.
  \item ESP32: complete the deadline loop (\texttt{TS\_US = 40}, \SI{25}{kHz})
    and verify that no sample is late. Analyse with
    \texttt{./alias\_scan esp.csv 25000 12}. Then set \texttt{TS\_US = 60}
    (\SI{16.67}{kHz}). Why is this particularly dangerous for a knock
    detector that integrates the energy in a band around \SI{6.5}{kHz}?
  \item Why must the sampling period be a multiple of \SI{20}{\micro s}?
    Consider the Uno at exactly \SI{8}{kHz} (\SI{125}{\micro s}): determine
    the sequence of time offsets between the ADC instants and the last chip
    update, its effect on the SNR (Exercise~4.1\,e) and a value of the chip
    attribute \texttt{fs\_khz} in \texttt{diagram.json} that fixes the
    problem.
  \item \texttt{./alias\_scan sim 5000 10} simulates the chip and the ADC
    (``virtual Wokwi'', also writes \texttt{alias\_sim.csv}). Compare its
    peaks with your Wokwi measurement.
\end{tasks}
\begin{hint}
The UART transmits about 11.5 characters per ms: printing a CSV line per
sample is impossible at kHz rates, therefore samples are buffered in RAM
(Uno: 512 samples = \SI{1}{KB} of its \SI{2}{KB}). Wokwi is not
cycle-accurate for the ESP32 -- if \texttt{late} is not 0, increase
\texttt{TS\_US} to the next multiple of 20.
\end{hint}
\end{exercise}

\begin{solution}
a) $\texttt{OCR1A}=2\cdot\texttt{TS\_US}-1=399$ (\SI{2}{MHz} timer clock).
\inputcode[firstline=32,lastline=53]{solutions/wokwi/sheet05_alias_uno/sketch.ino}
The Timer0 overflow ISR (\texttt{millis()}, every \SI{1.024}{ms}, a few
\si{\micro s} long) can delay the Timer1 ISR, i.e.\ cause sampling jitter
(Sheet~4); the Timer1 ISR itself starts with a constant latency.

b) Results of the virtual Wokwi (\texttt{sim}, 400\,ms, 39 knock and 38 valve
bursts); your Wokwi recording should agree within the resolution:
\begin{center}\small
\begin{tabular}{@{}lcccc@{}}\toprule
 & \multicolumn{2}{c}{predicted (5.1)} & \multicolumn{2}{c}{measured peaks}\\
$f_s$ [kHz] & knock modes & valve & knock bursts & valve bursts\\\midrule
5 (Uno)       & 1.5 / 0.5      & 2.0   & 1.51, 0.255 (34\,\%)  & 1.78, 0.98 (83\,\%)\\
8.33 (Uno)    & 1.833 / 2.167  & 0.333 & 1.85                  & 0.367\\
25 (ESP32)    & 6.5 / 10.5     & 8.0   & 6.50, 10.525 (24\,\%) & 8.025\\
16.67 (ESP32) & 6.5 / 6.167    & 8.0   & 6.483                 & 8.183\\\bottomrule
\end{tabular}
\end{center}
(all values in kHz; percentages: power relative to the largest peak). At \SI{5}{kHz} the
\SI{6.5}{kHz} knock appears at \SI{1.5}{kHz}: a burst of $4\tau=\SI{3.2}{ms}$
(16 samples) shows about 5 periods of \SI{1.5}{kHz}, i.e.\ the ``wrong''
frequency is clearly visible in the time signal. The second peak (alias of
\SI{10.5}{kHz} at \SI{0.5}{kHz}) is shifted to \SI{255}{Hz} because the
segment of only \SI{3.2}{ms} resolves no better than $\approx1/\SI{3.2}{ms}\approx\SI{300}{Hz}$
and the combustion vibration adds low-frequency content. The valve burst
segments (\SI{2}{ms}, 10 samples) have $\approx\SI{500}{Hz}$ resolution; the
peak at \SI{1.78}{kHz} instead of \SI{2.0}{kHz} and the side peak result from
the overlap with the mirrored component at $f_s-\SI{2}{kHz}=\SI{3}{kHz}$. At
\SI{8.33}{kHz} the two knock aliases (\SI{333}{Hz} apart) merge into one peak
at \SI{1.85}{kHz}, and the \SI{8}{kHz} valve burst looks like a slow
\SI{0.33}{kHz} transient. The Uno at \SI{8.33}{kHz} works only because the
CPU does nothing but sampling during the block ($\approx\SI{112}{\micro s}$ of
every \SI{120}{\micro s}); D0 must show regular \SI{120}{\micro s} steps.

c) Deadline loop:
\inputcode[firstline=45,lastline=55]{solutions/wokwi/sheet05_alias_esp32/sketch.ino}
At \SI{25}{kHz} all components lie below $f_s/2=\SI{12.5}{kHz}$ and appear at
their true frequencies (\SI{6500}{Hz}, \SI{10525}{Hz} with power ratio
$0.24\approx0.5^2$ as in the model, valve \SI{8025}{Hz}). At \SI{16.67}{kHz}
mode 2 aliases to \SI{6.17}{kHz} and merges with mode 1 into a single peak at
\SI{6.48}{kHz}. A band-pass detector at \SI{6.5}{kHz} integrates the energy of
both modes, and every disturbance between \SI{10}{kHz} and \SI{11}{kHz}
(e.g.\ ignition or injector noise) would be counted as knock -- aliasing
produces false knock detections that no digital filter can remove afterwards.

d) Sampling a held (staircase) signal with a period that is a multiple of
\SI{20}{\micro s} is equivalent to exact decimation of the chip's
\SI{50}{kHz} sequence. With \SI{125}{\micro s} the offset between ADC instant
and last chip update cycles through $125n \bmod 20 = 0, 5, 10, 15, 0,\dots$
\si{\micro s}. The mean (\SI{7.5}{\micro s}) is a harmless delay, the varying
part has $\sigma=\SI{5.6}{\micro s}$, i.e.\
$\mathrm{SNR}=-20\log_{10}(2\pi\cdot6500\cdot5.6\cdot10^{-6})\approx\SI{12.8}{dB}$
for mode 1 (\SI{8.7}{dB} for mode 2), concentrated in spurious lines at
$\pm f_s/4$ around each component because the offset is periodic. Fix:
\texttt{"fs\_khz": "40"} (update every \SI{25}{\micro s}, $125=5\cdot25$; the
chip's own Nyquist limit \SI{20}{kHz} is still above \SI{10.5}{kHz}) or
\texttt{"200"} (\SI{5}{\micro s}, more simulation load).

e) The simulation produces the peaks of the table in b); the Wokwi
recordings differ only by the random knock amplitudes/phases (the chip uses
its own random sequence) and the number of bursts.
\end{solution}

%-------------------------------------------------------------------------
\begin{exercise}{Reconstruction and anti-aliasing}{\JSLinux}
\textbf{Reconstruction.} \codefile{jslinux/sheet05/recon.c} represents a
``continuous'' sine on a grid of $M=40$ points per sampling interval
($f_s=\SI{25}{kHz}$, \num{600} samples) and computes the normalised error
$\mathrm{NMSE}=\sum e^2/\sum x^2$ of a reconstruction on this grid.
\begin{tasks}
  \item Implement zero-order hold, linear interpolation and sinc
    interpolation with $\pm K=16$ terms, once plain (truncated) and once with
    a Hann taper $w(u)=\tfrac12+\tfrac12\cos(\pi u/(K+1))$. Show analytically
    that the ZOH of a sine has
    $\mathrm{NMSE}=2\bigl(1-\sin(\omega T_s)/(\omega T_s)\bigr)$ and compare.
    Discuss the table (method vs.\ $f/f_s$).
  \item A \SI{15}{kHz} tone is sampled at \SI{25}{kHz} and sinc-reconstructed.
    What comes out? Predict frequency \emph{and} phase.
\end{tasks}
\textbf{Anti-aliasing.} \codefile{jslinux/sheet05/antialias.c} simulates the
analog world on a \SI{1}{MHz} grid; the AA filter (cut-off $f_c=\SI{12}{kHz}$)
is simulated on this grid, the ADC takes every $(\SI{1}{MHz}/f_s)$-th value.
Test signals: knock tones \SI{6.5}{kHz}/\SI{10.5}{kHz}; a disturbance at
$f_d=f_s-\SI{6.5}{kHz}$ that aliases exactly onto the knock frequency; white
noise (bandwidth \SI{500}{kHz}).
\begin{tasks}[resume]
  \item Implement the RC low-pass (exact update for an input held constant
    over $\Delta t$), the 4th-order Butterworth as a cascade of two
    2nd-order sections ($Q_k=1/(2\sin\frac{(2k-1)\pi}{2n})$, given biquad
    simulator) and the amplitude measurement at a given frequency in the
    sampled signal. Run the program for $f_s=25$ and \SI{50}{kHz} and filter
    orders 0, 1, 2, 4. Compare the alias amplitudes with $|H(f_d)|$.
  \item Which Butterworth order would be required for \SI{40}{dB} alias
    suppression in the band \SIrange{0}{10.5}{kHz} with $f_c=\SI{12}{kHz}$ at
    $f_s=\SI{25}{kHz}$ and at \SI{50}{kHz}? Formulate a design
    recommendation for the knock ADC.
\end{tasks}
\end{exercise}

\begin{solution}
a) Implementations:
\inputcode[firstline=38,lastline=72]{solutions/jslinux/sheet05/recon.c}
ZOH error of $x=\sin(\omega t+\varphi)$ at $t=nT_s+u$:
$e=\sin(\alpha+\omega u)-\sin\alpha$; averaged over the phase
$\alpha$: $\overline{e^2}=\tfrac12|\e^{\jj\omega u}-1|^2=1-\cos\omega u$;
averaged over $u\in[0,T_s)$: $1-\sin(\omega T_s)/(\omega T_s)$; division by the
signal power $1/2$ gives the formula. Program output (NMSE in dB):
\begin{center}\small
\begin{tabular}{@{}rrrrrrr@{}}\toprule
$f$ [Hz] & $f/f_s$ & ZOH & ZOH theory & linear & sinc & Hann-sinc\\\midrule
500   & 0.02 & $-22.98$ & $-22.79$ & $-56.83$ & $-42.94$ & $-97.22$\\
1000  & 0.04 & $-16.94$ & $-16.78$ & $-44.79$ & $-44.98$ & $-99.46$\\
2500  & 0.10 & $-9.06$  & $-8.89$  & $-28.94$ & $-47.59$ & $-101.24$\\
5000  & 0.20 & $-3.28$  & $-3.13$  & $-17.12$ & $-37.44$ & $-84.43$\\
8000  & 0.32 & $0.27$   & $0.41$   & $-9.43$  & $-32.66$ & $-70.23$\\
10500 & 0.42 & $2.02$   & $2.13$   & $-5.28$  & $-30.83$ & $-57.78$\\
12000 & 0.48 & $2.76$   & $2.83$   & $-3.38$  & $-27.58$ & $-11.36$\\\bottomrule
\end{tabular}
\end{center}
The ZOH agrees with the theory within \SI{0.2}{dB} (the fine grid contains
$u=0$ with zero error, hence slightly smaller values). The ZOH is a poor
reconstruction: even at $f=f_s/25$ the error is only \SI{17}{dB} below the
signal -- mainly due to the delay $T_s/2$ and the staircase; above
$f\approx f_s/3.5$ the error is larger than the signal. Linear interpolation
gains about \SI{12}{dB} per octave (error $\propto(\omega T_s)^4$) but is
useless near $f_s/2$. The truncated sinc has a frequency-independent floor of
about \SI{-30}{}\dots\SI{-48}{dB} caused by cutting the slowly decaying ($1/u$) kernel;
at low frequencies it is even worse than linear interpolation. The Hann taper
removes this floor ($-97$ to \SI{-101}{dB}), but its transition band costs
accuracy close to $f_s/2$ ($-11$\,dB at $0.48f_s$): the closer a signal
gets to Nyquist, the longer the interpolation kernel has to be.

b) $\sin(2\pi\,15000\,nT_s+0.7)=\sin(2\pi(15000-25000)nT_s+0.7)
=\sin(-2\pi\,10000\,nT_s+0.7)$: the samples are those of a \SI{10}{kHz} sine
with inverted phase, $-\sin(2\pi\,10\,\text{kHz}\,t-0.7)$. The program finds
a zero-crossing frequency of \SI{10000}{Hz} and NMSE $=\SI{3.00}{dB}$ against
the \SI{15}{kHz} original, but $\SI{-59.7}{dB}$ against this alias. The
reconstruction is perfect -- of the wrong signal.

c) Key code:
\inputcode[firstline=49,lastline=57]{solutions/jslinux/sheet05/antialias.c}
\inputcode[firstline=91,lastline=94]{solutions/jslinux/sheet05/antialias.c}
\inputcode[firstline=104,lastline=109]{solutions/jslinux/sheet05/antialias.c}
Results ($f_c=\SI{12}{kHz}$; ``noise'' = noise PSD in \SIrange{5}{11}{kHz}
after sampling relative to the unaliased analog noise PSD):
\begin{center}\small
\begin{tabular}{@{}cc|cc|cc|c@{}}\toprule
$f_s$ & order & gain 6.5k & gain 10.5k & alias of $f_d$ & $|H(f_d)|$ & noise\\\midrule
25 kHz & 0 & 0.00 & 0.00 & 0.0 dB & 0.0 dB & $+15.7$ dB\\
($f_d=18.5$ kHz) & 1 & $-1.12$ & $-2.47$ & $-5.3$ dB & $-5.3$ dB & $+2.0$ dB\\
 & 2 & $-0.36$ & $-2.01$ & $-8.2$ dB & $-8.2$ dB & $0.0$ dB\\
 & 4 & $-0.03$ & $-1.29$ & $-15.2$ dB & $-15.2$ dB & $0.0$ dB\\\midrule
50 kHz & 0 & 0.00 & 0.00 & 0.0 dB & 0.0 dB & $+12.7$ dB\\
($f_d=43.5$ kHz) & 1 & $-1.12$ & $-2.47$ & $-11.5$ dB & $-11.5$ dB & $-0.3$ dB\\
 & 2 & $-0.36$ & $-2.01$ & $-22.4$ dB & $-22.4$ dB & $-1.0$ dB\\
 & 4 & $-0.03$ & $-1.29$ & $-44.8$ dB & $-44.7$ dB & $-0.2$ dB\\\bottomrule
\end{tabular}
\end{center}
(gains in dB.) The sampled alias amplitude equals $|H(f_d)|$ exactly: the
filter acts before sampling, the sampler only moves the attenuated component
to \SI{6.5}{kHz}. Without filter, all noise of the \SI{500}{kHz} band folds
into $[0,f_s/2]$: the in-band noise rises by $10\log_{10}(500/12.5)=\SI{16}{dB}$
(measured \SI{15.7}{dB}) at \SI{25}{kHz} and by \SI{13}{dB} at \SI{50}{kHz}.
Already a first-order filter removes most of this broadband penalty, but at
$f_s=\SI{25}{kHz}$ even the 4th-order filter suppresses the disturbance by only
\SI{15}{dB} while it costs \SI{1.3}{dB} at \SI{10.5}{kHz}.

d) The first frequency that folds into \SIrange{0}{10.5}{kHz} is
$f_s-\SI{10.5}{kHz}$. Requirement $(f/f_c)^{2n}\ge10^4$:
$f_s=\SI{25}{kHz}$: $f/f_c=14.5/12=1.21$, $n\ge 2/\log_{10}1.21=24.3$, i.e.\
order 25 -- impossible. $f_s=\SI{50}{kHz}$: $f/f_c=39.5/12=3.29$,
$n\ge2/\log_{10}3.29=3.9$, i.e.\ order 4 (the table confirms
\SI{44.8}{dB} at \SI{43.5}{kHz}). Recommendation: oversample (e.g.\
\SIrange{50}{100}{kHz}), use a simple 2nd--4th-order analog AA filter
(Sallen-Key), and reduce the rate digitally afterwards (digital low-pass +
decimation, Sheet~10), where steep filters are cheap and exactly
reproducible.
\end{solution}

\begin{deliverables}
5.1: alias table, folding diagram, AA and band-pass sampling calculation.
5.2: completed sketches, CSV files of the Uno (two rates) and the ESP32 (two
rates), output of \texttt{alias\_scan} with comparison to 5.1, answers to c)
and d). 5.3: completed programs, both tables with discussion, filter-order
recommendation.
\end{deliverables}
